\documentclass[11pt]{article}
\usepackage[utf8]{inputenc}
\usepackage[T1]{fontenc}
\usepackage{amsmath,amssymb,amsthm}
\usepackage[margin=1.1in]{geometry}
\usepackage{hyperref}
\usepackage{booktabs}
\usepackage{pgfplots}
\usepgfplotslibrary{groupplots}
\pgfplotsset{compat=1.17}

\newtheorem{theorem}{Theorem}
\newtheorem{proposition}[theorem]{Proposition}
\newtheorem{lemma}[theorem]{Lemma}
\newtheorem{corollary}[theorem]{Corollary}
\theoremstyle{definition}
\newtheorem{remark}[theorem]{Remark}

\title{A waist at the binary collision: two-letter syzygy orbits\\ of the strong-force three-body problem}
\author{Oleksii Stavriianov\thanks{Independent researcher. ORCID \href{https://orcid.org/0009-0002-8504-3626}{0009-0002-8504-3626}, \texttt{oleksii.stavriianov@gmail.com}}}
\date{Version 1, November 2026}

\begin{document}
\maketitle

\begin{abstract}
Consider the planar three-body problem with the strong-force potential $U=\sum m_im_j/r_{ij}^2$ at zero energy and
zero angular momentum. Modulo symmetries its solutions are the geodesics of Montgomery's Jacobi-Maupertuis metric on
the shape sphere minus the three binary collisions, a pair of pants whose ends are asymptotic to cylinders. For equal
masses the two-letter syzygy classes $\dots ijij\dots$, the loops around a single binary collision, carry no
solution. We show that the Gaussian curvature at a binary collision is, to leading order, $K=\kappa_2\rho^2$ with
\[
  \kappa_2=-\frac{(m_i+m_j)\bigl[m_k^2(m_i+m_j)^3-m_i^2m_j^2M\bigr]}{4M\,m_i^4m_j^4},\qquad M=m_1+m_2+m_3,
\]
the leading coefficient of Montgomery's curvature formula at the collision. When $\kappa_2>0$, that is
$m_k^2(m_i+m_j)^3<m_i^2m_j^2M$, the end has a \emph{waist}, a loop shorter than the asymptotic cylinder, and the class
$(ij)$ is realised by a collision-free solution, periodic up to rotation, whose action lies below the cylinder's
circumference. At most one end can have a waist, that of the two heaviest bodies; for masses $(1,1,m)$ the condition
is $m<m^*=(1+\sqrt{65})/16\approx0.5664$. We compute these orbits at nine mass triples (all hyperbolic) and show that
the class is born from the cylinder at $m^*$: the action of the minimiser satisfies
$2\pi-S\ge35.12\,(m^*-m)^2(1+o(1))$, and the computed actions approach this rate. We give a short
Gauss-Bonnet proof that no two-letter orbit exists for equal masses, and, for masses without a waist at any end, reduce
the non-existence to three conditions on families of curves; we prove the first by interval arithmetic at five mass
triples and on the family $(1,1,m)$, $0.567\le m\le1.07$, and check the other two numerically at five triples. A map
of the whole mass triangle agrees with the waist criterion at all $741$ grid points. The orbits, the search, the threshold and the converse checks come with independent verification scripts.
\end{abstract}

\section{Introduction}

Montgomery \cite{M05} showed that for the planar three-body problem with an attractive $1/r^2$ potential, zero energy
and zero angular momentum, the solutions are, after reduction by rotations, translations and scaling, the geodesics of
a complete Riemannian metric on the shape sphere with the three binary collision points removed: a pair of pants whose
three ends are asymptotic to cylinders. For equal masses this metric has negative curvature except at the two Lagrange
points; as a consequence every free homotopy class of loops contains at most one closed geodesic, the syzygy sequences
give a complete symbolic dynamics, and the classes coded by two-letter sequences $\dots ijij\dots$, the loops around a
single binary collision, are not realised by any bounded solution \cite[Thm.~3]{M05}. Montgomery also gave the
curvature for arbitrary masses in closed form \cite[(7.12)-(7.13)]{M05} and showed that for most unequal masses it
takes both signs near the Lagrange points \cite[Thm.~6]{M05}. (We use the numbering of the arXiv version of \cite{M05}.)

The two-letter classes are the untied classes in the sense of Montgomery \cite{Mon98} (after Gordon
\cite{Gor75,Gor77}). In the realisation results for the strong force they are the exceptions: every other class is
realised, and the two-letter classes are stated not to be, without a restriction on the masses \cite[\S 2.1]{M05},
\cite[\S 1]{MM15}. Near each end the curvature is negative whenever $\kappa_2<0$ (Lemma~\ref{lem:K}), in particular for
equal masses, and the end then approaches its cylinder monotonically \cite[remark after (3.18)]{M05}; the proof of
non-existence for equal masses \cite[\S 6.3]{M05} uses this monotone approach together with the isosceles reflections,
which are isometries for equal masses. For the Newtonian problem, \cite[open question 6]{MM15} names planetary and
earth-moon-sun type configurations as realising the two-letter classes. Problem 4 of the list \cite{ACS12}
(Montgomery's) asks which syzygy sequences are realised in the Newtonian problem; its comment records, for the strong
force with equal masses, the correspondence of the bounded collision-free orbits with the sequences in which no letter
repeats and no pair of letters repeats forever \cite{M05}. With a waist the periodic sequence $\dots ijij\dots$ is
realised (Theorem~\ref{thm:exist}), so that description fails whenever an end has a waist; whether it holds when no
end has one is open (Section~\ref{sec:converse}).

The present note studies the two-letter classes for unequal masses. The curvature tends to zero at each end; the
question is the sign of the next term. We compute its leading coefficient $\kappa_2$ in closed form
(Lemma~\ref{lem:K}). Its sign is the limit at the collision of Montgomery's sign function \cite[(7.13)]{M05}, a fact
the supplement checks symbolically, and it has a direct geometric meaning: $\kappa_2>0$, i.e.\
$m_k^2(m_i+m_j)^3<m_i^2m_j^2M$, exactly when the end has a waist, a loop shorter than its asymptotic cylinder, so that
the approach of the end to its cylinder is not monotone. In that case the class $(ij)$ is realised
(Theorem~\ref{thm:exist}). For equal masses $\kappa_2=-5/6$, in accordance with \cite{M05}. We describe the orbits
numerically, show how they are born from the cylinder as the masses cross the threshold, and treat the converse:
without a waist, the non-existence reduces to three conditions on families of curves; we prove the first by
interval arithmetic at five mass triples and on the family $(1,1,m)$, $0.567\le m\le1.07$, and check the other two
numerically at five triples. Numerical studies of the zero-angular-momentum strong-force problem, mostly for equal
masses, include \cite{Moo93,DPS17}; the curvature of the same metric, and the cylinder at a collision for arbitrary
masses, are studied in \cite[\S 3.3.1]{KS16}. Computer-assisted proofs now settle questions about periodic three-body
orbits \cite{CGHLM24}. For the Newtonian problem with unequal masses large catalogues of
periodic orbits have been computed \cite{LJL18}; the strong-force problem offers a setting in which part of such a
picture can be proved.

\section{Setting}

Let $q_1,q_2,q_3\in\mathbb R^2$ with masses $m_i>0$, centre of mass at the origin, $U=\sum_{i<j}m_im_j/r_{ij}^2$,
energy $E=0$ and angular momentum $L=0$. Since $U$ is homogeneous of degree $-2$, the Lagrange-Jacobi identity gives
$\ddot I=4E=0$ for $I=\sum m_i|q_i|^2$, so a bounded solution has constant $I$, which we normalise to $1$.

\paragraph{Shape sphere.} With mass-weighted Jacobi vectors
$z_1=\sqrt{\mu_{12}}\,(q_2-q_1)$ and $z_2=\sqrt{\mu_3}\,\bigl(q_3-\tfrac{m_1q_1+m_2q_2}{m_1+m_2}\bigr)$, where
$\mu_{12}=m_1m_2/(m_1+m_2)$ and $\mu_3=m_3(m_1+m_2)/M$, we have $I=|z_1|^2+|z_2|^2$, and the Hopf map
$w=(|z_1|^2-|z_2|^2,\,2z_1\bar z_2)\in S^2$ identifies shapes (triangles up to rotation and scale) with points of the
unit sphere. The kinetic metric descends to the round sphere of radius $1/2$, $g_0=\tfrac14 g_{S^2}$. Each $r_{ij}^2$ is
an affine function of $w$; for instance $r_{12}^2=(1+w_1)/(2\mu_{12})$. The equator $w_3=0$ consists of the collinear
shapes; it is cut by the three collision points $C_{12},C_{13},C_{23}$ into three arcs, arc $k$ being the shapes with
body $k$ in the middle.

\paragraph{The metric.} Solutions with $E=0$, $L=0$, $I\equiv1$ project to geodesics of $g=2U\,g_0$ on
$P=S^2\setminus\{C_{12},C_{13},C_{23}\}$, and conversely \cite{M05}. The action $\int 2T\,dt$ of a solution equals
the $g$-length of its projection; $g$ is twice the metric of \cite[(3.9a)]{M05}, which is why lengths are
actions, and our $\rho$ below is the unit-sphere distance. Near $C_{ij}$ the metric is asymptotic to a Euclidean cylinder of circumference
$c_{ij}=2\pi\sqrt{2\mu_{ij}m_im_j}$, $\mu_{ij}=m_im_j/(m_i+m_j)$ \cite[\S 3]{M05}; Lemma~\ref{lem:E}
below gives the precise form. A closed geodesic gives a solution that
is periodic up to a rotation of the whole configuration.

\paragraph{Syzygies.} The syzygy sequence of a curve lists which body is in the middle at each crossing of the
equator. Cancelling repeated letters (stutters) reduces it; free homotopy classes of loops in $P$ correspond to
reduced periodic sequences up to cyclic shift, the reduced sequence of a closed curve being that of its class
\cite{M05}. The two-letter class $(ij)$, the loop around $C_{ij}$, has the sequence $\dots ijij\dots$.

\section{The ends and the curvature at a binary collision}

Fix the collision $C_{12}$; the other two are treated alike. Let $\rho$ be the unit-sphere distance from $C_{12}$ and
$\chi$ the angle about it. At $I=1$, $r_{12}^2=\sin^2(\rho/2)/\mu_{12}$, so
\[
  U=\frac{\mu_{12}m_1m_2}{\sin^2(\rho/2)}+V,
\]
where $V$, the potential of the pairs $13$ and $23$, is smooth near $C_{12}$ with value $V_0=m_3^2(m_1+m_2)^2/M$ there:
at the collision configuration body $3$ is at distance $D$ from the coincident pair with $\mu_3D^2=1$, so
$V_0=(m_1+m_2)m_3/D^2$. We write $R^2=2\mu_{12}m_1m_2$, so that $c_{12}=2\pi R$.

\begin{lemma}[the end in conformal coordinates]\label{lem:E}
With $u=-\log\tan(\rho/2)$,
\[
  g=R^2e^{2\psi}\,(du^2+d\chi^2),\qquad e^{2\psi}=\cos^2(\rho/2)\,(1+\varepsilon),\qquad
  \varepsilon=\frac{\sin^2(\rho/2)\,V}{\mu_{12}m_1m_2}.
\]
The function $\psi$ is smooth in $\zeta=\tan(\rho/2)e^{i\chi}$ near $\zeta=0$, $\psi=|\zeta|^2h(\zeta)$ with $h$ smooth and
$h(0)=\tfrac12\bigl(V_0/(\mu_{12}m_1m_2)-1\bigr)$. In particular $\psi$ and all its derivatives in $(u,\chi)$ are
$O(e^{-2u})$, uniformly in $\chi$.
\end{lemma}

\begin{proof}
$g_0=\tfrac14(d\rho^2+\sin^2\rho\,d\chi^2)=\tfrac14\sin^2\rho\,(du^2+d\chi^2)$ since $du=-d\rho/\sin\rho$, and
$\tfrac12U\sin^2\rho=2\mu_{12}m_1m_2\cos^2(\rho/2)(1+\varepsilon)$. The coordinate $\zeta$ is the stereographic coordinate
centred at $C_{12}$, $\sin^2(\rho/2)=|\zeta|^2/(1+|\zeta|^2)$ and $\cos^2(\rho/2)=1/(1+|\zeta|^2)$, so
$2\psi=-\log(1+|\zeta|^2)+\log(1+\varepsilon)$ is $|\zeta|^2$ times a smooth function with the stated value at $0$. Since
$|\zeta|=e^{-u}$, $\partial_u=-(\xi\partial_\xi+\eta\partial_\eta)$ and $\partial_\chi=\xi\partial_\eta-\eta\partial_\xi$
($\zeta=\xi+i\eta$), and both operators map $|\zeta|^2\cdot(\text{smooth})$ into itself.
\end{proof}

Thus each end is a half-cylinder of circumference $c_{12}$ up to an exponentially small conformal factor; in particular
$g$ is complete. Two consequences are used below. A loop in $\{u\ge u_0\}$ winding once around the end has length at
least $c_{12}\,e^{-Ce^{-2u_0}}$, and an arc joining the levels $u_1<u_2$ has length at least
$R\,e^{-Ce^{-2u_1}}(u_2-u_1)$, where $|\psi|\le Ce^{-2u}$ as in the lemma. The loops $u=\Lambda$ have total geodesic curvature $\oint\partial_u\psi\,d\chi=O(e^{-2\Lambda})$.
Finally, the reflection $w_3\mapsto-w_3$ is an isometry of $g$ for all masses (each $r_{ij}^2$ depends on $w_1,w_2$
only), so the three arcs of the equator are geodesics.

\begin{lemma}\label{lem:K}
The Gaussian curvature of $g$ satisfies
\begin{align*}
  K&=\kappa_2\,\rho^2+O(\rho^3),\\
  \kappa_2&=\frac{1}{4\mu_{12}m_1m_2}\Bigl(1-\frac{V_0}{\mu_{12}m_1m_2}\Bigr)
         =-\frac{(m_1+m_2)\bigl[m_3^2(m_1+m_2)^3-m_1^2m_2^2M\bigr]}{4M\,m_1^4m_2^4},
\end{align*}
independently of the direction.
\end{lemma}

\begin{proof}
By Lemma~\ref{lem:E}, $K=-R^{-2}e^{-2\psi}(\psi_{uu}+\psi_{\chi\chi})$ and $\psi=h(0)e^{-2u}+O(e^{-3u})$ with derivatives,
so $K=-4h(0)R^{-2}e^{-2u}+O(e^{-3u})$. Since $e^{-2u}=\tan^2(\rho/2)=\rho^2/4+O(\rho^4)$, this is
$K=-h(0)\rho^2/R^2+O(\rho^3)$, which is the first expression; the second follows from
$\mu_{12}m_1m_2=m_1^2m_2^2/(m_1+m_2)$ and the value of $V_0$. The direction does not enter because the leading term
of $V$ at $C_{12}$ is the constant $V_0$: the third body is equidistant from the coincident pair.
\end{proof}

\begin{remark}
A second derivation works on the unit sphere. Writing $g=e^{2\phi}g_{S^2}$ with $e^{2\phi}=U/2$, one has
$K=e^{-2\phi}(1-\Delta_S\phi)=(4-2\Delta_S\log U)/(2U)$, and $\Delta_S\log(1-\cos\rho)=-1$ exactly; expanding $\log U$ gives
the same $\kappa_2$. The $O(\rho^3)$ term comes from $2\sin\rho\,\partial_\rho V$ and depends on the direction. At equal
masses $\kappa_2=-5/6$. Exact rational expansions at ten mass triples, and numerical checks at the five triples of
Section~\ref{sec:converse} and its control, agree with the formula (supplement).
\end{remark}

\begin{corollary}[waist criterion]\label{cor:waist}
The conditions (a) $m_3^2(m_1+m_2)^3<m_1^2m_2^2M$, (b) $V_0<\mu_{12}m_1m_2$, (c) $\kappa_2>0$ are equivalent. The round
loops $\rho=\mathrm{const}$ about $C_{12}$ have length
\[
  \ell(\rho)=c_{12}\cos(\rho/2)\Bigl\langle\sqrt{1+\varepsilon}\Bigr\rangle_\chi
  =c_{12}\Bigl[1+\frac{\rho^2}{8}\Bigl(\frac{V_0}{\mu_{12}m_1m_2}-1\Bigr)\Bigr]+O(\rho^3),
\]
with $\langle\cdot\rangle_\chi$ the mean over $\chi$. Hence if $\kappa_2>0$ they are shorter than the cuff for small $\rho$
(the end has a \emph{waist}), and if $\kappa_2<0$ they are longer. For masses $(1,1,m)$ condition (a) reads $8m^2-m-2<0$,
i.e.\ $m<m^*=(1+\sqrt{65})/16$.
\end{corollary}

\begin{proof}
The equivalence of (a), (b), (c) is algebra. By Lemma~\ref{lem:E} the length is $\oint Re^{\psi}d\chi$, and
$\varepsilon=\rho^2V_0/(4\mu_{12}m_1m_2)+O(\rho^3)$; the term of $V$ linear in the direction contributes only at
order $\rho^3$. Numerically, the configuration-space loop lengths agree with this expansion at fourteen mass triples,
including $m=0.566$ and $0.567$ on the two sides of $m^*$ (supplement).
\end{proof}

At $\kappa_2=0$ ($m=m^*$ for masses $(1,1,m)$) the leading order decides nothing; we leave that single case aside.

\begin{proposition}\label{prop:one}
At most one end has a waist, and it is the binary of the two heaviest bodies.
\end{proposition}

\begin{proof}
Waists at the ends $12$ and $13$ would give $m_3^2(m_1+m_2)^3<m_1^2m_2^2M$ and $m_2^2(m_1+m_3)^3<m_1^2m_3^2M$, hence
$(m_1+m_2)^3(m_1+m_3)^3<m_1^4M^2$. But $(m_1+m_2)(m_1+m_3)=m_1M+m_2m_3>m_1M$, so $m_1^3M^3<m_1^4M^2$, i.e.\ $M<m_1$,
impossible. If the end $ij$ has a waist and $m_k\ge m_i$, write $s=m_i+m_j$; then $m_k^2s^3<m_i^2m_j^2(s+m_k)$, and
$(m_i/m_k)^2(s+m_k)\le s+m_i$ for $m_i\le m_k$, so $s^3<m_j^2(s+m_i)$; but $s^3-m_j^2(s+m_i)=m_i(m_i^2+3m_im_j+m_j^2)>0$.
So $m_k<m_i$, and likewise $m_k<m_j$.
\end{proof}

\section{Existence}

\begin{theorem}\label{thm:exist}
If the end $ij$ has a waist, the two-letter class $(ij)$ contains a closed geodesic of $g$ of length
$L^*=\inf\{\text{length of loops in }(ij)\}<c_{ij}$. Hence there is a collision-free solution with $E=0$, $L=0$, $I=1$,
periodic up to a rotation, realising the class $(ij)$, with action $L^*<c_{ij}$.
\end{theorem}

\begin{proof}
By the waist, $L^*<c_{ij}$. Take loops $\gamma_n$ in the class, parametrised proportionally to arc length, with lengths
$L_n\downarrow L^*$ and $L_n<L_1<c_{ij}$. Let $u_{ij},u_{ik},u_{jk}$ be the coordinates of Lemma~\ref{lem:E} at the three
ends. No loop of the class lies in an end $\{u_{ik}\ge u_0\}$: such an end is an annulus whose loops are powers of the
class $(ik)$, and the three boundary classes of the pair of pants are distinct and primitive. So each $\gamma_n$ meets
$\{u_{ik}\le u_0\}$, and by the bound on arcs it reaches at most depth $u_0+O(L_1)$ in that end; likewise in the end
$jk$. In the end $ij$, a loop of the class in $\{u_{ij}\ge\Lambda\}$ winds once and has length at least
$c_{ij}e^{-Ce^{-2\Lambda}}>L_1$ for $\Lambda$ large, so each $\gamma_n$ meets $\{u_{ij}\le\Lambda\}$ and again goes at most a bounded
depth further. Hence all $\gamma_n$ lie in one compact set $Q\subset P$. They are equi-Lipschitz, so by Arzel\`a-Ascoli
a subsequence converges uniformly to a loop $\gamma\subset Q$ of length at most $L^*$; once $\gamma_n$ is uniformly
closer to $\gamma$ than the injectivity radius of $Q$, the two are freely homotopic, so $\gamma$ is in the class and has
length $L^*$. A shortest loop in a non-trivial free homotopy class is locally length minimising, hence a smooth
closed geodesic.

For the solution: with $r=\sqrt I$ the kinetic metric on the configurations with $L=0$, modulo rotations, is
$dr^2+r^2g_0$, and the Jacobi-Maupertuis metric at $E=0$ is $2(U/r^2)(dr^2+r^2g_0)=2U\,(d\log r)^2+g$ with $U$ the
shape potential. It does not depend on $\log r$, so at zero conjugate momentum ($\dot I=0$) its geodesics are the curves
with $r$ constant whose shape is a $g$-geodesic. After the time change $dt=ds/(2U)$ these are solutions with $E=0$ and
$I=1$, and the horizontal lift through the Hopf map gives $L=0$; the lift of a closed curve closes up to a rotation. The
solution is collision-free because $\gamma$ avoids the three binary collisions and $I=1$ excludes the triple
collision. Its action $\int2T\,dt$ is the $g$-length $L^*$.
\end{proof}

\section{The orbits}\label{sec:orbits}

We computed the orbits by a one-parameter reversible search: a syzygy with body $1$ between bodies $2$ and $3$, all
velocities perpendicular to the line, constrained by $P=0$, $L=0$, $E=0$, with the position of body $1$ on the segment
as the parameter; a root of the along-line shape velocity at the next syzygy (body $2$ in the middle) closes the orbit
by reversibility after twice that time. Table~\ref{tab:orbits} lists the orbits found; for $m=0.57$, $0.58$, $0.59$,
$0.60$, $0.65$, $0.75$ (masses $(1,1,m)$) and for $(1,0.5,0.45)$, $(1,0.2,0.19)$ and the controls $(1,0.3,0.5)$,
$(0.5,0.3,1)$ (no waist at the end $12$; both have one at the end $13$ of their two heaviest bodies, as
Proposition~\ref{prop:one} requires) no root exists, consistent with Corollary~\ref{cor:waist} and
Proposition~\ref{prop:one}; all nineteen cases are rerun in the supplement. The search sees only orbits with this
reversing symmetry; orbits without it are the
subject of Section~\ref{sec:converse}. The class $(12)$ is primitive, so each orbit runs once around $C_{12}$ and its
syzygy word per period is $12$. Times are for $I=1$; the action is scale invariant.

\begin{table}[h]
\centering
\begin{tabular}{lcccccc}
\toprule
masses & action $S$ & cuff $c_{12}$ & period & trace & rotation \\
\midrule
$(1,1,0.40)$ & $5.6342554918$ & $2\pi$ & $1.71728$ & $53.50$ & $+146.0^\circ$ \\
$(1,1,0.50)$ & $6.1597496557$ & $2\pi$ & $1.40783$ & $8.102$ & $-116.9^\circ$ \\
$(1,1,0.52)$ & $6.2197385980$ & $2\pi$ & $1.18199$ & $4.988$ & $-89.0^\circ$ \\
$(1,1,0.53)$ & $6.2429769247$ & $2\pi$ & $1.02557$ & $3.872$ & $-73.3^\circ$ \\
$(1,1,0.54)$ & $6.2613331277$ & $2\pi$ & $0.82939$ & $3.017$ & $-56.1^\circ$ \\
$(1,1,0.55)$ & $6.2744359523$ & $2\pi$ & $0.58067$ & $2.413$ & $-37.0^\circ$ \\
$(1,1,0.56)$ & $6.2817961704$ & $2\pi$ & $0.25915$ & $2.068$ & $-15.5^\circ$ \\
$(1,0.5,0.30)$ & $3.4385697805$ & $2\pi/\sqrt3$ & $2.88203$ & $22.01$ & $-171.1^\circ$ \\
$(1,0.2,0.15)$ & $1.5489069692$ & $2\pi/\sqrt{15}$ & $6.38024$ & $19.94$ & $-169.0^\circ$ \\
\bottomrule
\end{tabular}
\caption{Closed orbits in the class $(12)$; syzygy word $12$ per period, shape closure $7\cdot10^{-12}$ to
$2\cdot10^{-9}$. Trace: of the linearised return map of the closed geodesic; rotation: of the whole configuration per
period, counterclockwise positive, in $(-180^\circ,180^\circ]$ (so $+146.0^\circ$ is the same as $-214.0^\circ$).}
\label{tab:orbits}
\end{table}

\paragraph{Stability.} The normal Jacobi equation $J''+KJ=0$ along each closed geodesic gives the $2\times2$
linearised return map of the geodesic flow on the shape sphere, i.e.\ of the reduced flow; the full relative periodic
orbit has in addition the trivial multipliers of the symmetries, the integrals and the dilation. The trace
(Table~\ref{tab:orbits}) exceeds $2$ for all nine orbits, so they are hyperbolic as closed geodesics of $g$, with multiplier
$\lambda=(\mathrm{tr}+\sqrt{\mathrm{tr}^2-4})/2$ from $53.5$ at $m=0.40$ down to $1.30$ at $m=0.56$, tending to $1$ as
$m\to m^*$. This is the behaviour of a non-degenerate closed geodesic of index $0$ on an orientable surface
\cite{Kli78}, although $K$ is positive on part of every orbit (its maximum on the orbit is $0.003$ at $m=0.56$ and
$0.92$ at $(1,0.2,0.15)$). The rotation per period is the geometric phase of the shape curve \cite{Mon96}: half the
solid angle it encloses on the unit shape sphere, which the supplement checks against the rotation to $0.05^\circ$.

\paragraph{Birth from the cylinder.} For masses $(1,1,m)$ the numerical threshold lies in $(0.56,0.57)$, containing
$m^*=0.56639$. The round loops have length $\ell(\rho)=c_{12}\bigl[1-A_2\rho^2+A_4\rho^4+O(\rho^6)\bigr]$, the odd orders
averaging out, with $A_2=(1-v)/8=\mu_{12}m_1m_2\,\kappa_2/2$, $v=V_0/(\mu_{12}m_1m_2)$, and
\[
  A_4=\frac{(1-v)(1+3v)}{384}+\frac{m_3^3(m_1+m_2)^3}{32M^2m_1^2m_2^2}.
\]
Indeed $\ell=c_{12}\cos(\rho/2)\langle\sqrt{1+\varepsilon}\rangle_\chi$ (Corollary~\ref{cor:waist}), the mean of $V$
over the circle is $V_0+\tfrac{\rho^2}{4}\Delta_SV(C_{12})+O(\rho^4)$, and since $r_{13}^2$ and $r_{23}^2$ are affine
in $w$ one finds $\Delta_SV(C_{12})=V_0+m_3\mu_3^2$; expanding to order $\rho^4$ gives the formula. On the threshold curve $v=1$, so
\[
  A_4=\frac{m_3^3(m_1+m_2)^3}{32M^2m_1^2m_2^2}>0
\]
for all masses: the second coefficient is positive along the whole boundary of the waist region. For masses $(1,1,m)$,
$A_4(m^*)=m^{*3}/(4(2+m^*)^2)=0.0068967$. Minimising over $\rho$, the shortest round loop has
\begin{align*}
  c_{12}-\min_\rho\ell&=\frac{c_{12}A_2^2}{4A_4}\bigl(1+o(1)\bigr)=C\,(m^*-m)^2\bigl(1+o(1)\bigr),\\
  C&=\frac{c_{12}A_2'(m^*)^2}{4A_4(m^*)}=\frac{\pi(4+m^*)^2}{32\,m^{*5}}=35.1205,
\end{align*}
at $\rho^2\approx A_2/(2A_4)\approx28.5\,(m^*-m)$. The minimiser of Theorem~\ref{thm:exist} has length
$L^*\le\min_\rho\ell$, so $2\pi-L^*\ge C\,(m^*-m)^2(1+o(1))$; that the computed symmetric orbits are the minimisers we
have not proved, but their actions satisfy $S<\min_\rho\ell$ at all nine masses (by $10^{-6}$ at $m=0.56$, supplement).
Numerically the ratios $(2\pi-S)/(m^*-m)^2$ from Table~\ref{tab:orbits} are $28.00$,
$29.48$, $30.36$, $31.37$, $32.57$, $34.01$ for $m=0.50$, $0.52$ to $0.56$, and those of the shortest round loops are
$31.9$, $33.4$, $34.2$, $34.7$, $34.9$ at $m^*-m=0.02$, $0.01$, $0.005$, $0.0025$, $0.00125$: both rise towards $C$
(Figure~\ref{fig:birth}). As $m\to m^*$ the orbit shrinks into the end (Figure~\ref{fig:shape}), with the period and the
binary separation tending to zero. The sign $A_4>0$ is what makes the shortest round loop leave the cuff continuously
at $m^*$; the search finds no orbit of the class above $m^*$ (no fold).

\paragraph{The motion.} At $L=0$ the binary axis and the line from the binary's centre of mass to body $3$ turn in
opposite senses (Figure~\ref{fig:plane}): over one period by $+344.5^\circ$ and $-15.5^\circ$ at $m=0.56$, and by
$+146.0^\circ$ and $-214.0^\circ$ at $m=0.40$. The reason is the angular momentum: a circular binary of the $1/r^2$
problem with zero energy has angular momentum exactly $\sqrt{2\mu_{12}m_1m_2}=R$, whatever its size, so at $L=0$ the
outer motion carries $-R$, and near the threshold the line to body $3$ turns back by about $R$ times the period
($14.8^\circ$ predicted, $15.5^\circ$ observed at $m=0.56$). Near the threshold the motion is hierarchical: a tight binary spins about
once per period while the third body swings slightly back (binary separation $0.29$, distance of body $3$ from the
binary's centre $1.48$ at $m=0.56$). Away from it the two distances are comparable ($1.05$ to $1.13$ against $1.05$ to
$1.16$ at $m=0.40$). The cylinder end corresponds to the circular binary motions of the $1/r^2$ two-body problem, which
have zero energy at every radius; the third body's potential $V_0$ competes with the factor $\cos(\rho/2)$ by which the
sphere $I=1$ shrinks the loops, and the waist appears when $V_0<\mu_{12}m_1m_2$. In terms of angular momentum: a pair
with reduced mass $\mu$ in the potential $k/r^2$ at zero energy reaches collision unless its angular momentum is at
least $\sqrt{2\mu k}$.
The binary on the cylinder sits exactly at its critical value $R$, and the outer pair, body $3$ against the coincident
binary ($\mu_3$, $k=m_3(m_1+m_2)$), has critical value $\sqrt{2\mu_3m_3(m_1+m_2)}=\sqrt{2V_0}$. At $L=0$ the outer motion
carries $-R$, and the criterion $V_0<\mu_{12}m_1m_2$ reads $\sqrt{2V_0}<R$: the waist exists exactly when the angular
momentum the binary forces on the outer motion exceeds the outer pair's own critical value. In Euler's two-centre problem the
corresponding orbits are the planetary orbits of \cite{DM16}, with fixed centres; for the Newtonian three-body problem
compare \cite[open question 6]{MM15}.

\begin{figure}[t]
\centering
\begin{tikzpicture}
\begin{axis}[width=0.5\textwidth, axis equal image, xmin=-1.6, xmax=1.6, ymin=-1.5, ymax=1.5,
  xlabel={$w_2/(1-w_1)$}, ylabel={$w_3/(1-w_1)$}, tick label style={font=\scriptsize},
  legend style={font=\scriptsize, at={(1.03,0.5)}, anchor=west}]
\addplot[gray, thin, forget plot] coordinates {(-1.6,0) (1.6,0)};
\addplot[black, thick] table {fig/shape_0.40.dat};
\addplot[blue, thick] table {fig/shape_0.50.dat};
\addplot[teal, thick] table {fig/shape_0.53.dat};
\addplot[red, thick] table {fig/shape_0.56.dat};
\addplot[only marks, mark=*, mark size=1.5pt] coordinates {(0,0)};
\legend{$m=0.40$, $m=0.50$, $m=0.53$, $m=0.56$, $C_{12}$}
\end{axis}
\end{tikzpicture}
\caption{The orbits of masses $(1,1,m)$ on the shape sphere, in the stereographic chart centred at $C_{12}$; the
horizontal axis is the equator (collinear shapes). As $m\to m^*$ the orbit contracts into the end.}
\label{fig:shape}
\end{figure}

\begin{figure}[t]
\centering
\begin{tikzpicture}
\begin{groupplot}[group style={group size=3 by 1, horizontal sep=0.9cm}, width=0.36\textwidth, axis equal image,
  tick label style={font=\scriptsize}, title style={font=\small}]
\nextgroupplot[title={$(1,1,0.40)$}]
\addplot[black, line width=1.6pt, opacity=0.45] table[x=x1, y=y1] {fig/plane_040.dat};
\addplot[blue, dashed] table[x=x2, y=y2] {fig/plane_040.dat};
\addplot[red] table[x=x3, y=y3] {fig/plane_040.dat};
\nextgroupplot[title={$(1,1,0.56)$}]
\addplot[black, line width=1.6pt, opacity=0.45] table[x=x1, y=y1] {fig/plane_056.dat};
\addplot[blue, dashed] table[x=x2, y=y2] {fig/plane_056.dat};
\addplot[red] table[x=x3, y=y3] {fig/plane_056.dat};
\nextgroupplot[title={$(1,0.2,0.15)$}]
\addplot[black, line width=1.6pt, opacity=0.45] table[x=x1, y=y1] {fig/plane_asym.dat};
\addplot[blue, dashed] table[x=x2, y=y2] {fig/plane_asym.dat};
\addplot[red] table[x=x3, y=y3] {fig/plane_asym.dat};
\end{groupplot}
\end{tikzpicture}
\caption{Three periods in the plane, centre of mass at the origin: body $1$ grey (wide), body $2$ blue (dashed),
body $3$ red; at $(1,1,0.56)$ the paths of bodies $1$ and $2$ nearly coincide. The
configuration returns rotated by the angle of Table~\ref{tab:orbits}.}
\label{fig:plane}
\end{figure}

\begin{figure}[t]
\centering
\begin{tikzpicture}
\begin{axis}[width=0.48\textwidth, xmode=log, xlabel={$m^*-m$}, ylabel={$(2\pi-S)/(m^*-m)^2$},
  tick label style={font=\scriptsize}, legend style={font=\scriptsize}, legend pos=south west, ymin=24, ymax=36.5,
  xmin=0.001, xmax=0.1]
\addplot[only marks, mark=*, mark size=1.5pt] table {fig/birth_ratio_orbits.dat};
\addplot[blue, mark=o, mark size=1.5pt] table {fig/birth_ratio_loops.dat};
\addplot[red, thick, domain=0.001:0.1] {35.12};
\legend{orbits, shortest round loops, $C=35.12$}
\end{axis}
\end{tikzpicture}
\hfill
\begin{tikzpicture}
\begin{axis}[width=0.46\textwidth, axis equal image, hide axis, xmin=-0.05, xmax=1.05, ymin=-0.08, ymax=0.92]
\addplot[black] table {fig/triangle.dat};
\addplot[fill=red!30, draw=red] table {fig/waist_12.dat};
\addplot[fill=red!30, draw=red] table {fig/waist_13.dat};
\addplot[fill=red!30, draw=red] table {fig/waist_23.dat};
\addplot[blue, dashed] table {fig/slice.dat};
\addplot[only marks, mark=*, mark size=1.3pt, blue] table {fig/mstar.dat};
\node[font=\scriptsize, anchor=north] at (axis cs:0,0) {$m_1$};
\node[font=\scriptsize, anchor=north] at (axis cs:1,0) {$m_2$};
\node[font=\scriptsize, anchor=south] at (axis cs:0.5,0.866) {$m_3$};
\end{axis}
\end{tikzpicture}
\caption{Left: the birth of the orbit from the cylinder, masses $(1,1,m)$: the ratio $(2\pi-S)/(m^*-m)^2$ for the
orbits of Table~\ref{tab:orbits} and for the shortest round loops about $C_{12}$, rising towards the constant $C$. Right: the mass triangle
($m_1+m_2+m_3=1$); shaded, the masses for which the end of the two bodies on the nearby edge has a waist
(Corollary~\ref{cor:waist}); the regions are disjoint (Proposition~\ref{prop:one}). Dashed, the masses $(1,1,m)$; dot,
$m=m^*$.}
\label{fig:birth}
\end{figure}

\section{The converse}

\subsection{Equal masses}

\begin{lemma}\label{lem:GB}
If $\gamma$ is a simple closed geodesic in the class $(ij)$, the region $A$ between $\gamma$ and the end $ij$ has
$\int_A K\,dA=0$.
\end{lemma}

\begin{proof}
A simple closed curve on the sphere minus three points separates the points as $0|3$ or $1|2$, so it is
null-homotopic or isotopic to the loops around one puncture; hence $\gamma$ is isotopic to the loops
$u_{ij}=\mathrm{const}$ and $A$ is an annulus running into the end. Cut it at $u_{ij}=\Lambda$ and apply
Gauss-Bonnet to the compact annulus (Euler characteristic $0$): $\gamma$ has zero geodesic curvature and the loop
$u_{ij}=\Lambda$ has total geodesic curvature $O(e^{-2\Lambda})$ by Lemma~\ref{lem:E}. The integral converges absolutely, as
$K\,dA=O(e^{-2u})\,du\,d\chi$ in the end.
\end{proof}

\begin{corollary}[{\cite[Thm.~3]{M05}}]
For equal masses no two-letter class is realised.
\end{corollary}

\begin{proof}
Here $K\le0$, with equality only at the two Lagrange points \cite{M05}. Let $\gamma$ be a closed geodesic in the class
$(ij)$, $\varphi$ the corresponding deck transformation of the universal cover $X$ of $(P,g)$, and
$\hat P=X/\langle\varphi\rangle$ the annular cover. $X$ is complete, simply connected and has $K\le0$, so its geodesics
are embedded lines and a lift $\tilde\gamma$ of $\gamma$ is an axis of $\varphi$ \cite[II.6]{BH99}. The preimage in $X$
of its image $\hat\gamma\subset\hat P$ is $\bigcup_k\varphi^k\tilde\gamma=\tilde\gamma$, so $\hat\gamma$ is a simple
closed geodesic in the core class of $\hat P$. The end $ij$ lifts isometrically to $\hat P$, its fundamental group being
$\langle\varphi\rangle$; $\hat\gamma$ cuts $\hat P$ into two annuli, one running into this end, and the proof of
Lemma~\ref{lem:GB} applies there word for word: $\int K\,dA=0$ over a region where $K\le0$, with $K<0$ off the discrete
preimage of the Lagrange points, which is impossible.
\end{proof}

\subsection{Unequal masses without a waist}\label{sec:converse}

Throughout this subsection no end has a waist: $\kappa_2<0$ at all three collisions, so by Lemma~\ref{lem:K} the
curvature is negative in a punctured neighbourhood of each collision. (By Proposition~\ref{prop:one} a waist can occur
at one end only; when it occurs at an end other than $12$, the class $(12)$ is not covered by what follows.) We reduce the
non-existence of a closed geodesic in the class $(12)$ to three conditions, each a check on one-parameter families of
curves.

\begin{lemma}[maximum principle]\label{lem:F}
Let $\tau$ be a $C^2$ function with $\nabla\tau\ne0$ on an open set $\Omega$, whose level sets have geodesic curvature
strictly positive towards increasing $\tau$; equivalently $\mathrm{Hess}\,\tau(v,v)<0$ for $v\ne0$ tangent to a level
set, the geodesic curvature being $-\mathrm{Hess}\,\tau(v,v)/|\nabla\tau|$ for unit $v$. If a
geodesic $\gamma$ is tangent at $\gamma(0)\in\Omega$ to the level set through it, then $\tau(\gamma(t))<\tau(\gamma(0))$
for all small $t\ne0$. In particular no geodesic arc in $\Omega$ has an interior local minimum of $\tau$, and no geodesic
touches a level set from the side of increasing $\tau$.
\end{lemma}

\begin{proof}
$(\tau\circ\gamma)'(0)=0$ and $(\tau\circ\gamma)''(0)=\mathrm{Hess}\,\tau(\gamma'(0),\gamma'(0))<0$.
\end{proof}

\paragraph{Conditions.} Let $\sigma_0$ be a smooth simple loop around $C_{12}$, bounding a disc $D\ni C_{12}$ that
contains no other collision, with geodesic curvature positive towards $C_{12}$ (convex), and disjoint from arc $3$.
Let $\Omega_0=P\setminus\overline D$.
\begin{itemize}
\item[(I)] $K\le0$ on $\Omega_0$.
\item[(II)] There are smooth convex loops $\sigma_0\supset\sigma_1\supset\dots\supset\sigma_N$ around $C_{12}$, each
inside the previous one, and times $T_k>0$, with the following property. For each $k<N$ the normal map
$\Phi_k(s,t)=\exp(t\,n(s))$, $s\in\sigma_k$, $0\le t\le T_k$, $n$ the inward unit normal, is injective; along every
normal the Jacobi field $G''=-KG$, $G(0)=1$, $G'(0)=-k_g(s)$ has $G>0$ and $G'<0$ on $[0,T_k]$ (the parallels are
immersed and convex); and the parallel $\Phi_k(\cdot,T_k)$ lies inside $\sigma_{k+1}$. The last loop $\sigma_N$ lies
in a cap $\rho<\rho_N$ about $C_{12}$ on which the round circles are convex towards $C_{12}$ (condition C:
$2\cot\rho+\partial_\rho\log U>0$).
\item[(III)] No stutter: no geodesic arc lies, apart from its two endpoints, in an open hemisphere and has both
endpoints on the same arc of the equator.
\end{itemize}
Near $C_{12}$ condition C holds automatically when $\kappa_2<0$: half its left side is the $\rho$-derivative of the
logarithm of the length density of the circle $\rho=\mathrm{const}$, which by Corollary~\ref{cor:waist} is
\[
  \tfrac12\bigl(2\cot\rho+\partial_\rho\log U\bigr)=\frac\rho4\Bigl(\frac{V_0}{\mu_{12}m_1m_2}-1\Bigr)+O(\rho^2).
\]

\begin{proposition}\label{prop:G}
Under (I), (II), (III) the class $(12)$ contains no closed geodesic.
\end{proposition}

\begin{proof}
Let $\gamma$ be a closed geodesic in the class $(12)$.

\emph{Syzygies.} The sequence of $\gamma$ is not empty, the class being non-trivial. A stutter in it is an arc as
excluded in (III), so the sequence is reduced; a reduced sequence determines the free homotopy class and conversely \cite{M05}, so it is $1212\dots$ and
$\gamma$ does not cross arc $3$. It does not touch it either, since arc $3$ is a geodesic.

\emph{$\gamma$ lies in $\overline D$.} The closure of $\Omega_0$ in $S^2$ is a disc containing $C_{13}$, $C_{23}$ and arc
$3$. A geodesic tangent to $\sigma_0$ stays locally outside $D$ (Lemma~\ref{lem:F} for the distance to $\sigma_0$), so
$\gamma$ crosses $\sigma_0$ transversally wherever it enters $D$, and by compactness meets $\sigma_0$ finitely
often. Suppose first that $\gamma$ meets $D$ and $\Omega_0$, and let $\beta$ be an arc of $\gamma$ in $\Omega_0$ with
endpoints on $\sigma_0$. If $\beta$ is not simple, its first return to one of its own points closes a geodesic monogon
in $\overline{\Omega_0}$, excluded as below. If $\beta$ is simple, it cuts the disc $\overline{\Omega_0}$ in two; arc $3$, which $\beta$
avoids, lies in one of them, and the other, $\Delta$, is a disc without punctures bounded by $\beta$ and an arc of
$\sigma_0$. Gauss-Bonnet on $\Delta$, with exterior angles $\theta_1,\theta_2\in(0,\pi)$ and the arc of $\sigma_0$ curving
away from $\Delta$, gives $\int_\Delta K=2\pi-\theta_1-\theta_2+\int_{\sigma_0\cap\partial\Delta}k_g>0$, against (I).
Suppose next that $\gamma\subset\overline{\Omega_0}$; it may touch $\sigma_0$ from outside at finitely many points.
Then $\gamma$ lies in the annulus $\overline{\Omega_0}\setminus$arc $3$, in its core class. If $\gamma$ is simple it
bounds with $\sigma_0$ an annulus $A'$, or, if it touches $\sigma_0$, a chain of discs with two cusps each (exterior
angles $\pi$); in both cases Gauss-Bonnet gives $\int K=\oint_{\sigma_0}k_g>0$ ($k_g$ towards $C_{12}$), against (I).
If not, $\gamma$ contains a geodesic monogon $\beta$ with interior angle $\varphi\in(0,\pi)$ (the same evaluations hold
if its corner lies on $\sigma_0$): if $\beta$ bounds a disc without punctures,
$\int K=\pi+\varphi>0$; otherwise $\beta$ is in the core class and bounds with $\sigma_0$ an annulus with
$\int K=\oint_{\sigma_0}k_g+\pi-\varphi>0$; both contradict (I).

\emph{The chain.} Let $k$ be the largest index with $\gamma$ in the closed disc bounded by $\sigma_k$. If $k=N$, let $p$
be a point of $\gamma$ of largest $\rho$; there $\gamma$ touches the circle $\rho=\rho(p)$ from the side of $C_{12}$,
which Lemma~\ref{lem:F} with $\tau=-\rho$ and condition C excludes. If $k<N$, some point $q$ of $\gamma$ is not inside
$\sigma_{k+1}$. The image of $\Phi_k$ is the closed annulus between $\sigma_k$ and the parallel at time $T_k$ (by
injectivity), and it contains every point inside $\sigma_k$ and not inside $\sigma_{k+1}$; on it $t=\Phi_k^{-1}$ is a
$C^2$ function with $|\nabla t|=1$ ($G>0$: no focal points) whose level sets are the parallels, convex towards increasing
$t$ ($\mathrm{Hess}\,t(v,v)=G'/G<0$ on unit tangents). Let $p$ minimise $t$ on the part of $\gamma$ in the annulus.
Then $t(p)\le t(q)<T_k$, so $\gamma$ stays in the annulus near $p$ and $p$ is a local minimum of $t\circ\gamma$; this
contradicts Lemma~\ref{lem:F}, in the interior if $t(p)>0$ and at the boundary $\sigma_k$ (where $\gamma$ would touch
$\sigma_k$ from inside) if $t(p)=0$, $\Phi_k$ extending to small negative $t$.
\end{proof}

\paragraph{Stutters.} Condition (III) reduces to the normal geodesics of the equator arcs. Let $H$ be an open hemisphere
and $j$ an equator arc. Suppose the normal geodesics $\nu_s$ of $j$ into $H$ (Fermi coordinates $g=dr^2+G^2ds^2$,
$G''=-KG$, $G(0)=1$, $G'(0)=0$) form a diffeomorphism of $\{(s,r):0<r<r_{\rm exit}(s)\}$ onto $H$, with $G>0$ and
$G'>0$. Then $r$ is $C^2$ on $H$, its level sets are convex towards decreasing $r$, and Lemma~\ref{lem:F} with $\tau=-r$
excludes an interior maximum of $r$ on a geodesic arc in $H$, hence any stutter arc on $j$. Independently, Gauss-Bonnet
limits the stutters that could occur: when the hemisphere's total positive curvature
$\Pi_H=\int_HK^+\,dA$ is below $\pi$ there is no geodesic monogon in $H$ ($\int K=\pi+\varphi$), so a stutter arc is
simple and bounds with the segment of $j$ between its endpoints (which contains no collision) a disc $B$ with
$\int_BK=\theta_p+\theta_q$; both corner angles are then at most $\Pi_H$. For equal masses $K\le0$ and this recovers the
absence of stutters \cite[Prop.~2.2]{KLM08}.

\paragraph{Numerical status.} At the five triples $(1,1,0.57)$, $(1,1,0.6)$, $(1,1,2)$, $(1,0.5,0.45)$,
$(1,0.2,0.19)$, none of which has a waist at any end, (I) is proved and (II), (III) hold numerically, with $\sigma_0$
and the chain taken from a curve-shortening flow started near the cuff.

(I) is proved by interval arithmetic. In the stereographic chart $\zeta$ centred at a collision $C_{ij}$ the metric is
$g=2U(1+|\zeta|^2)^{-2}|d\zeta|^2$, and with
$F=|\zeta|^2U(1+|\zeta|^2)^{-2}=m_im_j/(2a_{ij}(1+|\zeta|^2))+V|\zeta|^2(1+|\zeta|^2)^{-2}$, where
$r_{ij}^2=a_{ij}(1-c_{ij}\cdot w)$ and $V$ is the potential of the other two pairs, $F$ is rational, positive and smooth
at $\zeta=0$, and $K<0$ if and only if $Q=F\,\Delta F-|\nabla F|^2>0$ (as $\log|\zeta|^2$ is harmonic). We enclose $Q$
over boxes with second-order interval jets in outward-rounded IEEE arithmetic (only $+,-,\times,/$ and $\sqrt{}$, all
correctly rounded), the masses and all derived constants being intervals. The boxes cover, in the chart of each
collision, its Voronoi cell (the points no farther from it than from the other two), with adaptive refinement; a box is
set aside when it lies wholly outside the cell or wholly inside the polygon $\sigma_0$. At all five triples every
remaining box has $Q>0$ (between $1.7\cdot10^4$ and $9.3\cdot10^5$ boxes per triple), so $K<0$ on the whole region
outside $\sigma_0$, the collisions $C_{13}$, $C_{23}$ included; with a small loop in place of $\sigma_0$ the same
procedure fails, as it must. With the masses as wide intervals the same proof covers whole intervals of masses at
once: with the single loop $\sigma_0$ of $(1,1,0.57)$, $K<0$ outside $\sigma_0$ for every mass triple $(1,1,m_3)$ with
$0.567\le m_3\le1.07$ ($173$ mass intervals of width at most about $0.005$). On that range the loop is convex for the
metric of each mass, encloses $C_{12}$ only and avoids arc $3$ (checked numerically at $300$ masses), so it is a
loop $\sigma_0$ as in the conditions; for $m_3$ above about $1.077$ it is no longer convex. The test whether a box lies
inside $\sigma_0$ uses floating point with a margin. Numerically, the maximum of $K$ outside
$\sigma_0$, on a
$400\times800$ grid away from caps of radius $0.01$ about the collisions, is between $-6.1\cdot10^{-4}$ and
$-1.4\cdot10^{-5}$ (the last value at $(1,1,2)$ is attained at the edge of the caps about $C_{13}$ and $C_{23}$,
where $\kappa_2\rho^2=-1.3\cdot10^{-5}$; the caps $10^{-3}\le\rho\le0.1$ about all three collisions are checked
with (III) and have $K<0$, and below $10^{-3}$ Lemma~\ref{lem:K} applies), and $\sigma_0$ avoids arc $3$. (II): chains
of $261$, $153$, $43$, $35$ and $659$ loops, polygons of $200$ points with positive discrete curvature, the loops
saved whenever some point had moved a $g$-distance $0.01$ ($0.001$ at $(1,0.2,0.19)$, where the loops pass near the
very curved end $23$); for every pair one $T_k$, all $200$ normals keep $G'\le-2.9\cdot10^{-4}$ on $[0,T_k]$, the
parallels are simple and nested at $21$ times in $[0,T_k]$ (between the samples a pair of double points could only
appear and disappear through a tangency of two sheets), the parallel at $T_k$ lies inside the next loop, and
condition C holds out to the last loop. A sparser chain fails: the normals from the flat parts of a loop must then run
until the slowest normal arrives and lose their convexity, which is why the hypothesis is stated with one time per
loop. (III): on each of the three arcs, $160$ normals ($100$ from the arc outside caps of radius $0.05$ and $30$ from each
cap down to $\rho=10^{-3}$) keep $G\ge1$ and $G'>0$ and leave the hemisphere in monotone order; the normals from inside
a cap stay within $0.1$ of their collision, where $K<0$. The normals that run into the cap of another collision are
followed further: each either crosses the equator inside the cap or goes down the end (at $\rho<0.01$, moving inwards
at an angle too steep to turn back by the Clairaut relation of the near-cylinder end), and none returns to the
hemisphere; at $(1,1,0.57)$, for instance, $10$, $10$ and $28$ normals per arc do so, $3$, $3$ and $4$ of them crossing
the equator. Below $\rho=10^{-3}$ we rely on the near-cylinder form of Lemma~\ref{lem:E}, without a
quantitative estimate. $\Pi_H$ is between $3.8\cdot10^{-4}$ and $1.5\cdot10^{-2}$. With a waist, at $(1,1,0.5)$, the
same check fails on the two arcs ending at $C_{12}$ and passes on the third. The integrations are fixed-step
fourth-order Runge-Kutta in arclength, each repeated with half the step, and both step sizes must pass. Conditions
(II) and (III) are numerical evidence in double precision; making them rigorous, and covering all masses rather than
these triples, is left open.
The band just above $m^*$, where all margins vanish with $\kappa_2$, needs a separate argument.

\paragraph{The whole mass triangle.} Beyond these triples we mapped the triangle numerically
(Figure~\ref{fig:map}). At $741$ mass triples on a grid ($m_i=k_i/40$), we minimised the $g$-length over loops around
$C_{12}$, star-shaped in the chart about $C_{12}$ with the logarithm of the radius a trigonometric polynomial of degree
$12$, from three starting circles, with the other two collisions kept outside. By Theorem~\ref{thm:exist} a loop
shorter than the cuff already gives a closed geodesic. Such a loop is found at exactly the $198$ triples with a waist at
the end $12$, where the minimisation stops on the orbit; at the other $543$ every minimisation slides into the end with
length above the cuff, and no local minimum away from the end appears. The search sees minimisers and local minima
only, not closed geodesics of saddle type; the map comes from the research script \texttt{tools/words\_T0.py}
and has no independent checker.

\begin{figure}[t]
\centering
\begin{tikzpicture}
\begin{axis}[width=0.5\textwidth, axis equal image, hide axis, xmin=-0.05, xmax=1.05, ymin=-0.08, ymax=0.92]
\addplot[black] table {fig/triangle.dat};
\addplot[fill=red!15, draw=red] table {fig/waist_12.dat};
\addplot[only marks, mark=*, mark size=0.7pt, red] table {fig/T0_yes.dat};
\addplot[only marks, mark=*, mark size=0.5pt, gray] table {fig/T0_no.dat};
\node[font=\scriptsize, anchor=north] at (axis cs:0,0) {$m_1$};
\node[font=\scriptsize, anchor=north] at (axis cs:1,0) {$m_2$};
\node[font=\scriptsize, anchor=south] at (axis cs:0.5,0.866) {$m_3$};
\end{axis}
\end{tikzpicture}
\caption{The class $(12)$ over the mass triangle ($m_1+m_2+m_3=1$). Red dots: a loop shorter than the cuff was
found (the class is realised); grey dots: every length minimisation slid into the end. Shaded: the waist region
$m_3^2(m_1+m_2)^3<m_1^2m_2^2M$ of Corollary~\ref{cor:waist}.}
\label{fig:map}
\end{figure}

\section{Verification}

The supplement contains checkers that share no code with the research tools, each ending with ALL PASS:
\begin{itemize}
\item the initial conditions of the nine orbits of Table~\ref{tab:orbits} to full double precision, re-integrated with
a fixed-step Runge-Kutta scheme (two step sizes) to confirm the invariants, the shape closure, the action and the
syzygy word;
\item the reversible search of Section~\ref{sec:orbits} at nineteen mass triples, in a regularised time and at two
step sizes, with
every sign change of the shooting function refined to tell a root from a jump between branches: a root at the start
parameter of each of the nine orbits of Table~\ref{tab:orbits}, and none at the ten other triples;
\item a symbolic check that the curvature criterion is equivalent to the mass condition, gives $8m^2-m-2=0$ for masses
$(1,1,m)$, and is the limit at the collision of the sign function of \cite[(7.13)]{M05}; a configuration-space
computation of the loop lengths around $C_{12}$ on both sides of the threshold, which also confirms the size of
$\kappa_2$ through Corollary~\ref{cor:waist}; exact rational expansions of $K$ at a collision;
\item for condition (I) at the five triples and on the family $(1,1,m)$, $0.567\le m\le1.07$, the
interval-arithmetic proof described in Section~\ref{sec:converse}, with its negative control;
\item for the converse at the five triples of Section~\ref{sec:converse}, the chains of convex loops and a checker with
its own curvature formula, charts and integrator that re-checks Lemma~\ref{lem:K}, the absence of a waist at every
end, (I), (II) in the form stated, condition C, (III) with the end normals, and $\Pi_H$, together with the negative
control $(1,1,0.5)$.
\end{itemize}
Two further checkers recompute the traces and rotations of Table~\ref{tab:orbits} (the rotations also as half the
enclosed solid angle) and the birth law: $A_4(m^*)$, $C$, the ratios quoted in Section~\ref{sec:orbits}, and
$S<\min_\rho\ell$ at all nine masses.

The numerical searches, the symbolic and interval computations and the verification scripts were carried out with the
assistance of a large language model. All results are reproducible from the scripts and certificates supplied; the
checkers were written independently of the research code, so that they do not share the assumptions of the code that
produced the results.

\begin{thebibliography}{99}
\bibitem{ACS12} A.~Albouy, H.~E.~Cabral, A.~A.~Santos, Some problems on the classical $n$-body problem,
\emph{Celest. Mech. Dyn. Astr.} 113 (2012) 369-375, arXiv:1305.3191.
\bibitem{BH99} M.~R.~Bridson, A.~Haefliger, \emph{Metric Spaces of Non-positive Curvature}, Springer, 1999.
\bibitem{CGHLM24} R.~Calleja, C.~Garc\'ia-Azpeitia, O.~H\'enot, J.-P.~Lessard, J.~D.~Mireles James, From the Lagrange
triangle to the figure eight choreography: proof of Marchal's conjecture, arXiv:2406.17564 (2024).
\bibitem{DPS17} V.~Dmitra\v{s}inovi\'c, L.~V.~Petrovi\'c, M.~\v{S}uvakov, Periodic three-body orbits with vanishing
angular momentum in the Jacobi-Poincar\'e strong potential, \emph{J. Phys. A} 50 (2017) 435102, arXiv:1710.03687.
\bibitem{DM16} H.~R.~Dullin, R.~Montgomery, Syzygies in the two center problem, \emph{Nonlinearity} 29 (2016) 1212-1237,
arXiv:1509.05220.
\bibitem{Gor75} W.~B.~Gordon, Conservative dynamical systems involving strong forces, \emph{Trans. Amer. Math. Soc.}
204 (1975) 113-135.
\bibitem{Gor77} W.~B.~Gordon, A minimizing property of Keplerian orbits, \emph{Amer. J. Math.} 99 (1977) 961-971.
\bibitem{KLM08} S.~R.~Kaplan, M.~Levi, R.~Montgomery, Making the moon reverse its orbit, or, stuttering in the planar
three-body problem, \emph{Discrete Contin. Dyn. Syst. Ser. B} 10 (2008) 569-595.
\bibitem{Kli78} W.~Klingenberg, \emph{Lectures on Closed Geodesics}, Springer, 1978.
\bibitem{KS16} G.~S.~Krishnaswami, H.~Senapati, Curvature and geodesic instabilities in a geometrical approach to the
planar three-body problem, \emph{J. Math. Phys.} 57 (2016) 102901, arXiv:1606.05091.
\bibitem{LJL18} X.~Li, Y.~Jing, S.~Liao, Over a thousand new periodic orbits of a planar three-body system with unequal
masses, \emph{Publ. Astron. Soc. Japan} 70 (2018) 64, doi:10.1093/pasj/psy057, arXiv:1709.04775.
\bibitem{MM15} R.~Moeckel, R.~Montgomery, Realizing all reduced syzygy sequences in the planar three-body problem,
\emph{Nonlinearity} 28 (2015) 1919-1935, arXiv:1412.2263.
\bibitem{Mon96} R.~Montgomery, The geometric phase of the three-body problem, \emph{Nonlinearity} 9 (1996) 1341-1360,
arXiv:dg-ga/9510005.
\bibitem{Mon98} R.~Montgomery, The $N$-body problem, the braid group, and action-minimizing periodic solutions,
\emph{Nonlinearity} 11 (1998) 363-376.
\bibitem{M05} R.~Montgomery, Hyperbolic pants fit a three-body problem, \emph{Ergodic Theory Dynam. Systems} 25 (2005)
921-947; arXiv:math/0405014 (``Fitting hyperbolic pants to a three-body problem'').
\bibitem{Moo93} C.~Moore, Braids in classical dynamics, \emph{Phys. Rev. Lett.} 70 (1993) 3675-3679.
\end{thebibliography}

\end{document}
